\chapter{MACVIO: Learning-Based Visual-Inertial SLAM with Cross-Modal Uncertainty Fusion} \label{chapter:macvio}
\fancyhead[LO]{\makebox[\textwidth][r]{Chapter 8. MACVIO}}

\section{Introduction}

This chapter presents MACVIO (Metrics-Aware Cross-Modal Visual-Inertial Odometry), which extends the single-sensor uncertainty learning from MAC-VO and AirIMU to multi-sensor fusion scenarios. MACVIO demonstrates how cross-modal uncertainty correlations enable adaptive fusion weights and uncertainty-driven sensor selection.

\section{Cross-Modal Uncertainty Correlation Learning}

\subsection{Problem Formulation}

Traditional visual-inertial SLAM systems rely on fixed fusion weights or simple heuristics for sensor selection. This approach fails to capture the dynamic correlation between visual and inertial uncertainty estimates that varies across environments and motion patterns.

% TODO: Develop MACVIO problem formulation and methods
[Content to be developed based on MACVIO research]

\subsection{Learning Cross-Modal Dependencies}

Building on the uncertainty models established in MAC-VO (Chapter 3) and AirIMU (Chapter 4), this section develops methods for learning the correlation structure between visual and inertial uncertainties.

[Content to be developed]

\section{Adaptive Sensor Fusion Weights}

\subsection{Uncertainty-Driven Weight Adaptation}

[Content to be developed]

\subsection{Real-Time Fusion Strategy}

[Content to be developed]

\section{Uncertainty-Driven Sensor Selection}

\subsection{Selective Integration Framework}

[Content to be developed]

\subsection{Performance under Sensor Degradation}

[Content to be developed]

\section{Visual-Inertial SLAM with Learned Uncertainty}

\subsection{System Architecture}

[Content to be developed]

\subsection{Loop Closure with Uncertainty Estimates}

[Content to be developed]

\section{Experimental Validation}

\subsection{Benchmark Datasets}

[Content to be developed]

\subsection{Comparison with State-of-the-Art}

[Content to be developed]

\subsection{Cross-Platform Evaluation}

[Content to be developed]

\section{Conclusion}

[Content to be developed based on research outcomes]